While the model assumes the droplet is at least partially ionized, the plasma state is not formed without a laser source. The laser introduces a fourth particle -- the photon -- into the system, and it is quickly
absorbed by either bound or free electrons. Whereas the term ``electrons'' has typically referred to free electrons for much of this chapter, in this context it collectively refers to both free and bound electrons.

\subsection{Energy deposition}
This model assumes the laser beam is a line heat source within each droplet even if the laser cross-sectional area is known. The volumetric heat generation, as a result, contains a Dirac delta function
$\delta(\bm{r}_{\text{laser}})$ where $\bm{r}_{\text{laser}}$ is the laser trajectory. Finite-volume methods blend well with this assumption since the delta function is always integrated out during the formulation.
\begin{equation}
    \iiint_{\Gamma \supset \bm{r}_0 } \delta(\bm{r}_0) dV = 1
\end{equation}

Therefore, the total energy deposited in a cell that the laser passes through is the difference in laser energy at the cell entrance and exit. If the entrance energy upon cell $k$ is $I_k A$, it will have been
attenuated by a factor of $e^{-\alpha_k s_k}$ when exitting the cell -- $s_k$ being the distance that the laser beam travels within cell $k$. Therefore,

\begin{equation}
    \iiint_{V_k} q'''_{\text{laser}} dV = I_k A (1 - e^{-\alpha_k s_k})
\end{equation}

Bound and free electrons differ in how the energy is spent. A free electron, upon absorbing a photon, simply increases in kinetic energy. A bound electron, however, can either be excited or removed from its orbit altogether.
If excitation occurs, collisions with the atom causes the electron to rapidly relax and transfer the acquired energy into thermal energy of the bulk liquid; therefore, on the timescale of interest, it is responsible
to assume that the energy goes directly into heating the neutral fluid.

With two channels in which laser energy can be deposited, the total attenuation coefficient $\alpha$ is the sum of the bound and free counterparts.
\begin{equation}
    \alpha = \alpha_b + \alpha_{\text{ib}},
    \label{eqn:total_attenuation_coeff}
\end{equation}
and the fraction of energy that goes into each channel is proportional to the relative magnitude of the two absorption coefficients
\begin{subequations}
    \begin{align}
        q_b &= I_0 A(1 - e^{-\alpha s}) \frac{\alpha_b}{\alpha} \label{eqn:heat_bound_electrons} \\
        q_{\gamma e} &= I_0 A(1 - e^{-\alpha s}) \frac{\alpha_{\text{ib}}}{\alpha}.
    \end{align}
\end{subequations}

The subscript ``ib'' in Equation \ref{eqn:total_attenuation_coeff} refers to inverse bremsstrahlung, a process in which free electrons, in the vincinty of an ion, absorbs energy and starts to accelerate. The double
subscript $\gamma e$ infers as if the heat source were due to a ``collision'' between a photon and a free electron, similar to how other heat sources are denoted in Section \ref{sect:relaxations}. Inverse bremsstrahlung
should take over as the dominant absorption process once a sufficiently opaque medium is formed. The attenuation coefficient will increase proportionally to the inverse bremsstrahlung frequency $\nu_{\text{ib}}$
\begin{equation}
    \alpha_{\text{ib}} = \frac{\nu_{\text{ib}} \eta}{c} = \frac{\omega_e^2}{\omega^2 + \nu_{ei}^2} \frac{\nu_{ei} \eta}{c},
    \label{eqn:nu_ib_general}
\end{equation}
where $\omega_e$ is the plasma electron frequency and $\omega$ is the laser frequency. The former is defined as
\begin{equation}
    \omega_e^2 \equiv \frac{n_e e^2}{\varepsilon_0 m_e}.
\end{equation}

When $\omega \gg \nu_{ei}$, $\nu_{\text{ib}}$ is usually given in the form
\begin{equation}
    \nu_{\text{ib}} = \frac{n_e}{n_c} \nu_{ei},
    \label{eqn:nu_ib_approx}
\end{equation}
and the critical density $n_c$ is defined in terms of $\omega$
\begin{equation}
    \omega^2 = \frac{n_c e^2}{\varepsilon_0 m_e}.
\end{equation}

\subsection{Photoionization}
\label{sect:photoionization}
A bound electron requires a binding energy, $\mathcal{E}_b$, to remove it from the orbit. When a laser beam deposits energy into an atom, each photon contributes an energy of $h\nu$. Therefore, the minimum number of
photons needed to completely remove an electron from its orbit is
\begin{equation}
    N \equiv \left\lceil \frac{\mathcal{E}_b}{h\nu} \right\rceil.
\end{equation}

The $N$-photon ionization rate is proportional to the $N$-th power of a particular measure of photon flux. In Mihelic (2021), that measure is the photon scalar flux, $\displaystyle \phi = \frac{I}{h\nu}$ \cite{Mihelic}. However,
since the laser intensity monotonic decreases with distance travelled, it is better to calculate the ionization rate using a path-averaged scalar flux within a computational cell $k$:
\begin{subequations}
    \begin{align}
        \overline{\phi}_k &= \frac{\overline{I}_k}{h\nu} \\
        &= \frac{1}{h\nu s_k} \int_{0}^{s_k} I_0 e^{-\alpha_k s} ds \\
        &= \frac{I_0}{h\nu} \frac{1 - e^{-\alpha_k s_k}}{\alpha_k s_k}
    \end{align}
\end{subequations}

The rate is then given by
\begin{equation}
    \Gamma_{\text{pi}} = S_{\text{pi}}^{(N)} \overline{\phi}^N n_g
    \label{eqn:photoionization_rate}
\end{equation}

Here, $S_{\text{pi}}^{(N)}$ is the $N$-photon ionization rate coefficient and $n_g$ is the neutral particle number density in the gas phase. In order for $\Gamma_{\text{pi}}$ to have the correct unit, the unit of
$S_{\text{pi}}^{(N)}$ is m$^{2N}$ s$^{N-1}$. If $N = 1$, then $S_{\text{pi}}^{(1)}$ has the expected units of cross section or area. For that reason, $S_{\text{pi}}^{(N)}$ is often referred to as a cross section
in literature -- regardless of $N$ and how photon concentration is expressed (other than $\overline{\phi}$) -- despite the fact that it can assume different units. For consistency, the symbol $\sigma$ reserved for
cross section is only used to denote quantities with units of area.

In the case of single-photon ionization, $S_{\text{pi}}^{(1)}$ is related to the absorption coefficient $\alpha_b$ by
\begin{equation}
    S_{\text{pi}}^{(1)} = \Phi \frac{\alpha_b}{n_g},
    \label{eqn:S_pi1}
\end{equation}
where $\Phi \in [0, 1]$ is the quantum efficiency, which denotes the fraction of absorbed photons that actually induce ionization. If $\Phi = 1$, then $\alpha_b$ is equivalent to a macroscopic cross section.

Upon ionization, the electron fluid gains an amount of energy equivalent to
\begin{equation}
    q^e_{\text{pi}} = (Nh\nu - \mathcal{E}_b) \Gamma_{\text{pi}}
    \label{eqn:electron_energy_N},
\end{equation}
while any energy from the initial laser beam not spent on ionizing will instead serve as a heat source for the neutral atoms:
\begin{equation}
    q'''_{\gamma n} = q'''_b - Nh\nu \Gamma_{\text{pi}}
    \label{eqn:excitation_energy}.
\end{equation}

For simplicity, bound electrons are assumed to only absorb $N$ photons prior to removal from its parent nucleus. Above-threshold ionization, a phenomenon where an electron absorbs more than $N$ photons, is also possible,
but will be neglected here.